% Part: sets-functions-relations
% Chapter: arithmetization
% Section: rationals
% Hindi translation (hi-Deva-IN) of Open Logic commit 9620cc73f9c8e0ad003c514a5d3748f29611c4c0.

\documentclass[../../../include/open-logic-section]{subfiles}

\begin{document}

\olfileid[hi]{sfr}{arith}{rat}
\olsection{$\Int$ से $\Rat$ तक}

अभी हमने देखा कि सरल समुच्चय-सिद्धांत का उपयोग करके प्राकृत संख्याओं से
पूर्णांक कैसे बनाए जाते हैं। अब लगभग उसी प्रकार पूर्णांकों से परिमेय संख्याएँ
बनाएँगे। हमारे प्रारंभिक अवलोकन हैं:
\begin{enumerate}
	\item प्रत्येक परिमेय संख्या को $\nicefrac{i}{j}$ के रूप में लिखा जा सकता
	है, जहाँ $i$ और $j$ दोनों पूर्णांक हैं, किंतु $j$ शून्येतर है।
	\item व्यंजक $\nicefrac{i}{j}$ में निहित जानकारी को क्रमित युग्म
	$\tuple{ i, j}$ में भी समान रूप से रखा जा सकता है।
\end{enumerate}
स्वाभाविक उपाय परिमेय संख्याओं को
$\Int \times (\Int \setminus \{0_\Int\})$ से लिए गए क्रमित युग्म
\emph{मानना} होगा। पर पहले की तरह यह अत्यधिक सरल होगा, क्योंकि हम चाहते हैं कि
$\nicefrac{3}{2} = \nicefrac{6}{4}$ हो, जबकि $\tuple{ 3, 2}\neq \tuple{ 6,
4}$। अधिक सामान्य रूप में हमें यह चाहिए:
\[
	\nicefrac{a}{b} = \nicefrac{c}{d} \text{ तभी और केवल तभी जब } a \times d = b \times c
\]
इसके लिए $\Int \times
(\Int \setminus \{0_\Int\})$ पर एक तुल्यता संबंध इस प्रकार परिभाषित करें:
\[
	\tuple{ a, b }\Ratequiv \tuple{ c, d} \text{ तभी और केवल तभी जब }a \times d = b \times c
\]
हमें जाँचना है कि यह तुल्यता संबंध है। यह~$\Intequiv$ के मामले से बहुत मिलता
है, इसलिए इसे अभ्यास के रूप में छोड़ते हैं।
\begin{prob}
दिखाइए कि $\Ratequiv$ एक तुल्यता संबंध है।
\end{prob}
तब हम यह परिभाषा दे सकते हैं:
\begin{defn}
परिमेय संख्याएँ, पूर्णांकों के उन युग्मों के~$\Ratequiv$ के अधीन तुल्यता वर्ग
हैं जिनका दूसरा अवयव शून्येतर है। अर्थात $\Rat =
\equivclass{(\Int \times (\Int\setminus \{0_\Int\}))}{\Ratequiv}$।
\end{defn}

पूर्णांकों की तरह यहाँ भी कुछ मूल संक्रियाएँ परिभाषित करनी हैं। जहाँ
$\equivrep{i,j}{\Ratequiv}$ से~$\Ratequiv$ के अधीन वह तुल्यता वर्ग सूचित हो
जिसका !!a{element} $\tuple{i, j}$ है, वहाँ हम कहते हैं:
\begin{align*}
	\equivrep{a, b}{\Ratequiv} + \equivrep{c, d}{\Ratequiv} &= \equivrep{ad + bc,  bd}{\Ratequiv}\\
	\equivrep{a, b}{\Ratequiv} \times \equivrep{c, d}{\Ratequiv} &= \equivrep{a  c, b d}{\Ratequiv}.
\intertext{इन परिमेयों पर $r \leq s$ परिभाषित करने के लिए इस तथ्य का उपयोग
करें कि $r \le s$ तभी और केवल तभी जब $s - r$ ऋणात्मक न हो; अर्थात $s - r$
को $\nicefrac{i}{j}$ के रूप में लिखा जा सके, जहाँ $i$ ऋणेतर और $j$~धनात्मक हो:}
	\equivrep{a, b}{\Ratequiv} \leq \equivrep{c, d}{\Ratequiv} &\text{ तभी और केवल तभी जब }
	\equivrep{c, d}{\Ratequiv} - \equivrep{a, b}{\Ratequiv} = 
	\equivrep{i_\Int, j_\Int}{\Ratequiv}
\end{align*}
जहाँ कोई $i \in \Nat$ और $0 \neq j \in \Nat$ हो।

फिर जाँचना होगा कि ये परिभाषाएँ वैसा ही व्यवहार करती हैं जैसा उन्हें
\emph{करना चाहिए}; यह जाँच \olref[check]{sec} में दी गई है। और वे सचमुच ऐसा
करती हैं। अंततः पूर्णांकों को परिमेय संख्याएँ \emph{मानकर} बरतने का कोई ढंग
चाहिए; अतः प्रत्येक $i \in \Int$ के लिए नियत करें कि $i_\Rat =
\equivrep{i, 1_\Int}{\Ratequiv}$। इन सबके उचित व्यवहार की जाँच भी
\olref[check]{sec} में की गई है।

\begin{prob}
दिखाइए कि $(i + j)_\Rat = i_\Rat+ j_\Rat$, $(i \times j)_\Rat =
i_\Rat \times j_\Rat$ और $i \leq j \liff i_\Rat \leq j_\Rat$, किन्हीं
भी $i, j \in \Int$ के लिए।
\end{prob}

\end{document}
